\section{Experimental Results}
\label{sec:results}

\subsection{Held-Out Text-to-Spark-SQL Performance}

The main held-out evaluation compares the three configurations defined in
Section~\ref{sec:methodology} over 20 frozen natural-language questions.

Table~\ref{tab:heldout-abc} reports validation rate, execution rate, result
accuracy, and the absolute number of correct queries.

\begin{table}[htbp]
\centering
\caption{Held-out Text-to-Spark-SQL evaluation.}
\label{tab:heldout-abc}
\begin{tabular}{lrrr}
\toprule
Metric & A Full & B Retrieval & C Repair \\
\midrule
Validation rate & 75\% & 85\% & 90\% \\
Execution rate & 75\% & 85\% & 90\% \\
Result accuracy & 70\% & 80\% & 85\% \\
Correct queries & 14/20 & 16/20 & 17/20 \\
\bottomrule
\end{tabular}
\end{table}


Figure~\ref{fig:heldout-abc} presents the same comparison graphically.

\begin{figure}[htbp]
    \centering
    \includegraphics[
        width=0.82\textwidth
    ]{heldout_abc_rates.pdf}
    \caption{
        Held-out Text-to-Spark-SQL performance for configurations A, B, and C.
    }
    \label{fig:heldout-abc}
\end{figure}

The complete-catalog baseline, Configuration A, achieves a validation rate of
75\%, an execution rate of 75\%, and a final result accuracy of 70\%. This
corresponds to 14 correct results out of 20 held-out questions.

Replacing the complete metadata catalog with relation-aware SQL-safe metadata
retrieval improves all three metrics. Configuration B reaches 85\% validation,
85\% execution, and 80\% result accuracy, corresponding to 16 correct results
out of 20.

Configuration C adds validation-driven one-shot repair and obtains the highest
final performance: 90\% validation, 90\% execution, and 85\% result accuracy,
or 17 correct results out of 20.

The improvement from A to B is particularly important because both
configurations use the same language model and deterministic generation
parameters. The main experimental difference is therefore the metadata context
supplied before SQL generation.

The results indicate that reducing the catalog to a relation-aware SQL-safe
subset does not merely reduce prompt size. Under the held-out benchmark, it
also improves final query correctness.


\subsection{Prompt Efficiency}

The metadata-retrieval configuration substantially reduces the amount of
context supplied to the language model.

Table~\ref{tab:efficiency-latency} summarizes prompt-size and controlled
latency measurements.

\begin{table}[htbp]
\centering
\caption{Prompt-size and controlled-latency comparison.}
\label{tab:efficiency-latency}
\begin{tabular}{lrrr}
\toprule
Metric & A Full & B Retrieval & Reduction \\
\midrule
Mean prompt tokens & 2558.45 & 776.35 & 69.66\% \\
Prompt evaluation (s) & 182.58 & 54.24 & 70.29\% \\
Wall-clock time (s) & 202.31 & 69.29 & 65.75\% \\
\bottomrule
\end{tabular}
\end{table}


On the 20 held-out questions, Configuration A uses an average of 2558.45 prompt
tokens, while Configuration B uses 776.35 tokens. The relation-aware approach
therefore reduces mean prompt size by 69.66\%.

Figure~\ref{fig:prompt-tokens} highlights this reduction.

\begin{figure}[htbp]
    \centering
    \includegraphics[
        width=0.68\textwidth
    ]{prompt_token_efficiency.pdf}
    \caption{
        Mean prompt-token count for complete-catalog and relation-aware
        prompting on the held-out benchmark.
    }
    \label{fig:prompt-tokens}
\end{figure}

The reduction is achieved without changing the underlying LLM. It results from
selecting only metadata considered relevant to the current analytical request
and rendering that metadata using physical SQL-safe identifiers.

This result supports one of the main motivations of the project: as metadata
catalogs grow, sending the complete catalog for every natural-language request
is not necessary for the evaluated workload.


\subsection{Controlled Inference Latency}

The sequential held-out execution exhibited strong cache effects, especially
for the complete-catalog prompts that repeatedly share a large common prefix.
For this reason, the primary latency comparison uses the controlled
prefix-isolated benchmark described in Section~\ref{sec:methodology}.

The controlled experiment reports a mean prompt-evaluation time of
182.58 seconds for Configuration A and 54.24 seconds for Configuration B.
This corresponds to a 70.29\% reduction in prompt-evaluation time.

Mean wall-clock time decreases from 202.31 seconds to 69.29 seconds, a
65.75\% reduction.

Figure~\ref{fig:controlled-latency} shows the controlled prompt-evaluation
comparison.

\begin{figure}[htbp]
    \centering
    \includegraphics[
        width=0.68\textwidth
    ]{controlled_latency_prompt_eval.pdf}
    \caption{
        Mean controlled prompt-evaluation latency for complete-catalog and
        relation-aware prompting.
    }
    \label{fig:controlled-latency}
\end{figure}

These absolute latency values are specific to the CPU-only local execution
environment and should not be interpreted as general model-performance
measurements. The relevant result is the relative reduction observed when the
same model processes the smaller relation-aware prompt.

The experiment therefore links metadata-context reduction to a measurable
decrease in inference cost under controlled conditions.


\subsection{Metadata Retrieval Quality}

Before considering catalog growth, the relation-aware retriever was evaluated
on the frozen 20-question held-out benchmark using the original four-source
catalog.

The retriever achieves macro recall of 0.925, macro precision of 0.662, and
macro F1 of 0.754. Complete required-metadata recall is obtained for 17 of the
20 held-out questions.

The remaining three queries contain incomplete metadata retrieval. These
failures are especially informative because downstream Text-to-Spark-SQL
accuracy is strongly associated with metadata completeness.

Among the 17 perfect-recall questions, Configuration B returns the correct
result for 15 questions, corresponding to 88.2\% accuracy. Configuration C
returns the correct result for 16 of the 17, corresponding to 94.1\%.

For the three imperfect-recall questions, both B and C obtain only one correct
result, corresponding to 33.3\%.

This difference does not prove a causal relationship by itself, but it
provides strong evidence that missing required metadata is a major downstream
failure mode. Validation and repair can improve a query when the necessary
physical information is already present, but they cannot reliably reconstruct
relationships or schema elements that were absent from the prompt context.


\subsection{Repair Behaviour}

Configuration C triggered three repair attempts over the 20 held-out queries.
One of the three repairs produced a query that passed validation and improved
the final result, giving a validator-success rate of one third for attempted
repairs.

The successful case occurred for a Green Taxi geographic query. The initial
generation referenced an invalid pickup timestamp column. Spark validation
identified the schema error, and the one-shot repair removed the unnecessary
date restriction, producing a valid and correct query.

The other repair attempts corresponded to cases in which the retrieval context
itself lacked required metadata. In those cases, validator feedback alone was
insufficient to reconstruct the missing physical relationship.

A separate limitation also appears in valid-but-incorrect queries. For one
cross-taxi comparison, the SQL was executable and produced the correct
numerical aggregates but generated different categorical labels from those
required by the benchmark. Since validation checks syntactic and structural
validity rather than complete semantic equivalence, no repair was triggered.

The results therefore show that repair is useful as a limited recovery
mechanism for invalid SQL, but it is not a substitute for correct metadata
grounding or result-level evaluation.


\subsection{Catalog Scalability}

The catalog-scale experiment evaluates how retrieval behaves as the number of
metadata sources grows from 4 to 8 and 16.

Exact measurements are reported in Table~\ref{tab:catalog-scaling}.

\begin{table}[htbp]
\centering
\caption{Relation-aware retrieval under catalog scaling.}
\label{tab:catalog-scaling}
\begin{tabular}{rrrrrrrrr}
\toprule
Sources & Columns & Docs & Full chars & Mean ctx & Perfect & Recall & Precision & F1 \\
\midrule
4 & 66 & 103 & 9498 & 2160.7 & 17/20 & 0.925 & 0.662 & 0.754 \\
8 & 90 & 131 & 12024 & 2061.6 & 14/20 & 0.825 & 0.620 & 0.683 \\
16 & 139 & 188 & 16651 & 2061.6 & 14/20 & 0.825 & 0.620 & 0.683 \\
\bottomrule
\end{tabular}
\end{table}


The complete rendered catalog grows from 9,498 characters at four sources to
12,024 characters at eight sources and 16,651 characters at sixteen sources.

In contrast, the mean relation-aware context remains bounded. It decreases
slightly from 2160.7 characters at four sources to 2061.6 characters at eight
sources and remains at 2061.6 characters at sixteen sources.

Figure~\ref{fig:catalog-context} illustrates the difference between complete
catalog growth and retrieved-context size.

\begin{figure}[htbp]
    \centering
    \includegraphics[
        width=0.78\textwidth
    ]{catalog_scale_context_size.pdf}
    \caption{
        Complete catalog size and mean relation-aware context size as the
        catalog grows from 4 to 16 sources.
    }
    \label{fig:catalog-context}
\end{figure}

The corresponding context reduction increases from 77.25\% at four sources to
82.85\% at eight sources and 87.62\% at sixteen sources.

However, this increased compression is accompanied by a retrieval-quality
trade-off. Macro recall decreases from 0.925 at four sources to 0.825 at eight
sources. Macro precision decreases from 0.662 to 0.620, and macro F1 decreases
from 0.754 to 0.683.

The number of perfect-recall held-out questions decreases from 17 out of 20 to
14 out of 20.

Figure~\ref{fig:catalog-recall} shows the recall trend.

\begin{figure}[htbp]
    \centering
    \includegraphics[
        width=0.68\textwidth
    ]{catalog_scale_recall.pdf}
    \caption{
        Mean metadata recall under increasing catalog scale.
    }
    \label{fig:catalog-recall}
\end{figure}

Importantly, no additional degradation is observed when the catalog grows from
8 to 16 sources. Recall, precision, F1, mean context size, and perfect-recall
count remain unchanged between these two scales.

The result therefore does not indicate a simple monotonic relationship in
which every additional source causes further degradation. Instead, particular
semantically similar hard negatives introduced at the eight-source scale are
responsible for the observed loss.


\subsection{Catalog-Scale Error Analysis}

Post-hoc analysis of the frozen catalog-scale results identifies the specific
metadata lost when the first hard-negative sources are introduced.

For both held-out variants of the Green Taxi drop-off borough query, retrieval
is perfect with four sources. At eight sources, the retriever loses:

\begin{itemize}

    \item \texttt{taxi\_zones};
    \item \texttt{taxi\_zones.Borough};
    \item \texttt{taxi\_zones.LocationID};
    \item the Green Taxi drop-off-to-zone relationship.

\end{itemize}

A similar effect appears for one Yellow Taxi pickup-borough query, where the
retriever loses the Taxi Zones dataset, the borough and location columns, and
the Yellow pickup-to-zone relationship.

The most difficult three-source paraphrase already had incomplete weather
retrieval at four sources. After the hard-negative catalog expansion, it also
loses the geographic Taxi Zones metadata.

These failures are consistent with semantic competition from the added
domain-related distractors, particularly sources containing borough,
geographic, weather, or alternative mobility concepts.

The fact that the same failures persist unchanged at sixteen sources suggests
that the observed degradation is caused mainly by the first set of strongly
competitive hard negatives rather than by catalog size alone.


\subsection{Metadata Index Construction Cost}

Catalog growth also increases the amount of metadata that must be embedded and
indexed.

With four sources, the index contains 103 metadata documents and requires
6,472 embedding prompt tokens. Index construction uses 28.54 seconds of
embedding wall time.

At eight sources, the index contains 131 documents, requires 8,437 embedding
tokens, and uses 30.68 seconds of embedding wall time.

At sixteen sources, the index grows to 188 documents and 11,585 embedding
tokens, with 42.43 seconds of embedding wall time.

These measurements represent index-construction cost rather than per-query
retrieval cost. Since metadata indexes can be reused across many analytical
requests, this cost is operationally different from the prompt and generation
cost incurred for each user query.


\subsection{Spark Data-Volume Scalability}

The final experiment evaluates the physical execution layer independently from
metadata-catalog growth.

Table~\ref{tab:data-volume-scaling} reports median Spark runtime across three
repetitions for each workload and data volume.

\begin{table}[htbp]
\centering
\caption{Spark data-volume scaling using median runtime over three repetitions.}
\label{tab:data-volume-scaling}
\begin{tabular}{lrrrrr}
\toprule
Scale & Rows & Parquet MiB & Scan (s) & Grouping (s) & Weather join (s) \\
\midrule
100k & 100,000 & 2.31 & 0.140 & 0.226 & 0.274 \\
500k & 500,000 & 11.44 & 0.140 & 0.293 & 0.406 \\
1m & 1,000,000 & 23.24 & 0.146 & 0.224 & 0.615 \\
Full & 2,963,713 & 72.58 & 0.199 & 0.492 & 1.350 \\
\bottomrule
\end{tabular}
\end{table}


The physical Yellow Taxi benchmark scales from 100,000 rows to the complete
2,963,713-row curated dataset. The corresponding two-file Parquet
representation grows from 2.31 MiB to 72.58 MiB.

Figure~\ref{fig:data-volume} presents the runtime curves.

\begin{figure}[htbp]
    \centering
    \includegraphics[
        width=0.80\textwidth
    ]{spark_data_volume_scaling.pdf}
    \caption{
        Median Spark runtime under increasing Yellow Taxi data volume.
    }
    \label{fig:data-volume}
\end{figure}

For the simple scan-and-aggregation workload, median runtime is 0.140 seconds
at 100,000 rows and 0.199 seconds for the complete dataset. The corresponding
runtime factor is 1.42 despite an approximately 29.6-fold increase in row
count.

The pickup-zone grouping workload increases from 0.226 seconds to
0.492 seconds, corresponding to a 2.18-fold runtime increase.

The temporal taxi--weather join shows the clearest volume-dependent growth.
Median runtime increases from 0.274 seconds at 100,000 rows to 0.406 seconds
at 500,000 rows, 0.615 seconds at one million rows, and 1.350 seconds for the
complete dataset. The full-data runtime is 4.92 times the 100,000-row
baseline.

The first repetition of several small workloads is substantially slower than
later repetitions, confirming the presence of JVM, filesystem, and execution
warm-up effects. For this reason, the pre-defined benchmark reports the median
of three runs rather than interpreting a single execution.


\subsection{Interpretation of the Scaling Results}

The data-volume experiment should not be interpreted as evidence that Spark
has a universally sub-linear runtime curve. The maximum test dataset is only
72.58 MiB in its materialized Parquet representation and is executed using two
local worker threads.

At the smaller scales, fixed costs such as query planning, task scheduling,
JVM warm-up, and filesystem caching can dominate the actual scan cost. This is
particularly visible in the simple aggregation workload.

The temporal weather join provides a clearer scaling signal because it
performs additional timestamp transformation and cross-source matching.

The main conclusion from this experiment is therefore limited and
implementation-specific: the prototype successfully executes the selected
analytical workloads from 100,000 rows to the complete 2.96-million-row
Yellow Taxi dataset, while the measured runtime remains manageable in the
two-core experimental environment.


\subsection{Summary of Main Findings}

The experimental results provide four main findings.

First, relation-aware SQL-safe metadata retrieval improves held-out result
accuracy from 70\% to 80\% relative to complete-catalog prompting, and
validation-driven one-shot repair further increases accuracy to 85\%.

Second, the relation-aware context reduces mean prompt size by 69.66\% and
reduces controlled prompt-evaluation time by 70.29\% in the CPU-only
experimental environment.

Third, relation-aware retrieval keeps prompt context approximately bounded as
the metadata catalog grows, but semantically related hard negatives introduce
a measurable recall trade-off. Macro recall decreases from 0.925 to 0.825
when the first distractor sources are introduced, then remains stable as the
catalog grows further.

Fourth, the Spark execution layer successfully handles the evaluated
data-volume range up to 2,963,713 curated Yellow Taxi records, with the
temporal weather join exhibiting the clearest increase in runtime as volume
grows.

These findings are examined further in the discussion of system benefits,
limitations, and failure modes in Section~\ref{sec:discussion}.
